\section{Conceptos fundamentales}

\textbf{Espectrofotometría:} Técnica que mide la cantidad de luz absorbida o transmitida por una muestra en función de la longitud de onda.
\textbf{Absorbancia (A) y transmitancia (T):} Relacionadas por $A = -\log(T)$.


\section{Ley de Lambert-Beer}
La absorbancia es directamente proporcional a la concentración del analito:
\[
A = \varepsilon \cdot b \cdot c
\]
Donde:
\begin{itemize}
    \item $\varepsilon$: Absortividad molar (constante para una molécula y una longitud de onda dadas).
    \item $b$: Longitud del camino óptico (ancho de la celda).
    \item $c$: Concentración del analito.
\end{itemize}

\section{Intervalos del espectro}
\begin{itemize}
    \item \textbf{UV (Ultravioleta):} 200-400 nm. Se usa para compuestos con dobles enlaces conjugados.
    \item \textbf{Visible:} 400-700 nm. Se usa para compuestos coloreados.
    \item \textbf{IR (Infrarrojo):} >700 nm. Se usa para identificar grupos funcionales.
\end{itemize}

\section{Técnicas espectrofotométricas}
\begin{itemize}
    \item \textbf{UV-Visible:} Determinación de concentraciones de compuestos orgánicos e inorgánicos.
    \item \textbf{Espectrofotometría de absorción atómica (EAA):} Determinación de metales en concentraciones traza.
    \item \textbf{Infrarrojo (IR):} Identificación de compuestos por sus grupos funcionales.
\end{itemize}

\section{Curvas de calibración}
Se preparan estándares de concentraciones conocidas y se mide su absorbancia. La curva relaciona la absorbancia con la concentración y permite determinar la concentración de la muestra.

\resumebox{ejercicio}{Fuerza magnética del protón}{\footnotesize
    \textbf{Determinación de concentración por UV-Vis:} Se preparó una curva de calibración para la vitamina B12 con una ecuación $y = 0.02x + 0.01$ (donde $y$ es la absorbancia y $x$ es la concentración en $\mu$g/mL). Si la absorbancia de una muestra es de 0.41, ¿cuál es su concentración?
}

\resumebox{dato}{Fuerza magnética del protón}{\footnotesize
    \textbf{Control de calidad en la industria láctea:} La espectrofotometría UV-Vis se utiliza para determinar la concentración de vitaminas (ej. riboflavina), la detección de adulteraciones (ej. adición de agua), y la medición de la actividad de enzimas en la leche.
}